% 第五章 系统设计与实现
% ch5-system.tex

\subsection{系统功能概述}

\subsubsection{系统目标}

开发一个奥运项目评估决策支持系统，提供：
\begin{enumerate}
\item 数据管理与预处理功能
\item AHP-EWM混合权重计算
\item 项目综合评分与排名
\item 可视化展示与报告生成
\end{enumerate}

\subsubsection{系统架构}

\begin{figure}[htbp]
\centering
\begin{tikzpicture}[scale=0.8, every node/.style={scale=0.8}]
  % 数据层
  \node[draw, rounded corners, fill=blue!20, minimum width=8cm, minimum height=1.2cm] at (0,0) {数据层：数据存储、预处理、标准化};
  % 业务层
  \node[draw, rounded corners, fill=green!20, minimum width=8cm, minimum height=1.2cm] at (0,2) {业务层：AHP计算、EWM计算、混合权重};
  % 展示层
  \node[draw, rounded corners, fill=orange!20, minimum width=8cm, minimum height=1.2cm] at (0,4) {展示层：Web界面、可视化图表、报告生成};
  % 箭头
  \draw[->, thick] (0,0.7) -- (0,1.3);
  \draw[->, thick] (0,2.7) -- (0,3.3);
\end{tikzpicture}
\caption{系统架构图}
\label{fig:architecture}
\end{figure}

系统采用三层架构：
\begin{itemize}
\item 数据层：数据存储与预处理
\item 业务层：权重计算与评分逻辑
\item 展示层：Web界面与可视化
\end{itemize}

\subsection{核心模块实现}

\subsubsection{数据处理模块}

实现数据清洗、标准化、缺失值处理等功能。

\subsubsection{AHP权重计算模块}

实现判断矩阵构建、特征值计算、一致性检验。

\subsubsection{EWM权重计算模块}

实现信息熵计算、熵权求解。

\subsubsection{混合权重模块}

实现线性组合和灵敏度分析。

\subsection{可视化仪表盘}

\subsubsection{功能设计}

\begin{enumerate}
\item 项目评分雷达图
\item 权重分布柱状图
\item 历史趋势折线图
\item 预测结果排序图
\end{enumerate}

\subsubsection{技术选型}

\begin{itemize}
\item 后端：Python + Flask
\item 前端：Streamlit
\item 可视化：Plotly + ECharts
\end{itemize}

\subsection{系统演示}

\subsubsection{主界面}

系统主界面展示六维指标权重分布和项目综合评分。

\subsubsection{评估功能}

用户可输入项目指标数据，系统自动计算综合评分。

\subsubsection{预测功能}

基于历史数据和模型，预测未来奥运项目入选可能性。

\subsection{本章小结}

本章介绍了奥运项目评估系统的设计与实现。系统采用Python技术栈，实现了数据处理、权重计算、可视化展示等核心功能，为用户提供直观的评估决策支持。